\documentclass[11pt,a4paper]{article}

% ------------------------------------------------------------------ %
% Packages (mirrors the main solver report preamble)
% ------------------------------------------------------------------ %
\usepackage[margin=2.2cm]{geometry}
\usepackage{amsmath,amssymb,mathtools}
\usepackage{booktabs}
\usepackage{array}
\usepackage{xcolor}
\usepackage{listings}
\usepackage{siunitx}
\usepackage{microtype}
\usepackage{enumitem}
\usepackage{algorithm}
\usepackage{algpseudocode}
\usepackage[hidelinks]{hyperref}

\lstdefinestyle{cfg}{
  basicstyle=\ttfamily\footnotesize,
  breaklines=true,
  columns=fullflexible,
  frame=single,
  rulecolor=\color{black!30},
  showstringspaces=false,
}
\lstset{style=cfg}

% ------------------------------------------------------------------ %
% Convenience macros
% ------------------------------------------------------------------ %
\newcommand{\Vm}{V_{m}}
\newcommand{\Iion}{I_{\text{ion}}}
\newcommand{\Iext}{I_{\text{ext}}}
\newcommand{\Dten}{\mathbf{D}}
\newcommand{\grad}{\nabla}
\newcommand{\divg}{\nabla\!\cdot}
\newcommand{\code}[1]{\texttt{\small #1}}
\newcommand{\gpode}{\code{gaussPointODE}}
\newcommand{\ccrec}{\code{cellCentredReconstruct}}
\newcommand{\states}{\mathbf{s}}
\newcommand{\Vbar}{\overline{V}_{m}}
\newcommand{\Vbarc}{\overline{V}_{m,c}}
\newcommand{\avg}[1]{\left\langle #1 \right\rangle}
\newcommand{\Omegac}{\Omega_{c}}

\title{Why \ccrec\ states fail on a propagating TNNP front\\
       while \gpode\ works\\[2pt]
       \large A diagnosis of ignition and conduction failure in
       \texttt{highOrderElectroActivationFoamImplicitPETSc}}
\author{Solver diagnosis note}
\date{\today}

\begin{document}
\maketitle

\begin{abstract}
The solver offers two ways of handling the ionic state variables
$\states=(s_1,\dots,s_{N})$ in the high-order ionic-current quadrature:
\gpode\ (the legacy path, one stiff ODE integration per Gauss point) and
\ccrec\ (one ODE per cell centre, states then reconstructed at the Gauss points
by a high-order LRE of order \code{p1}/\code{p2}/\code{p3}). On the
Niederer\,\emph{et al.}~(2012) benchmark with the TNNP cell model, \ccrec\ fails
to produce a propagating activation as the time step is refined
($\Delta t=\SI{e-5}{s}$): the wavefront neither ignites (with the nominal
stimulus) nor conducts (even with a stronger stimulus), whereas \gpode\ ignites
and propagates on the same mesh and $\Delta t$. This note shows, with three
controlled discrimination runs and a short mathematical argument, that this is
\emph{not} a coding bug, nor a time-scheme (Crank--Nicolson) artefact, nor a
$\Delta t$-scaling error. It is an intrinsic accuracy limitation of evaluating a
strongly non-linear ionic kinetics at the \emph{cell-average} membrane potential:
by Jensen's inequality the cell-average sodium activation under-estimates the
pointwise average, weakening the regenerative $I_{\mathrm{Na}}$ current at the
sharp upstroke. Raising the reconstruction order (\code{p2}, \code{p3}) does not
help, because the order only affects how the already-computed states are
interpolated, not the potential that drove the ODE.
\end{abstract}

% ================================================================== %
\section{Governing equations and the two state-integration modes}
% ================================================================== %

The monodomain problem solved for the transmembrane potential $\Vm$ is
\begin{equation}
  \chi C_m \,\frac{\partial \Vm}{\partial t}
  \;=\; \divg\!\left(\Dten\,\grad \Vm\right)
        \;-\; \Iion(\Vm,\states) \;+\; \Iext ,
  \qquad
  \frac{\partial \states}{\partial t} \;=\; \mathbf{g}(\Vm,\states),
  \label{eq:monodomain}
\end{equation}
where $\states$ collects the $N$ gating variables and ion concentrations of the
ten Tusscher--Noble--Noble--Panfilov (TNNP) cell model, $\mathbf{g}$ is its stiff
right-hand side, and $\Iion=\phi(\Vm,\states)$ is the (algebraic) total ionic
current. The reaction source of \eqref{eq:monodomain} is integrated over each
cell $\Omegac$ with a high-order quadrature at points $\{x_q\}$ and weights
$\{w_q\}$ ($\sum_q w_q=|\Omegac|$),
\begin{equation}
  \overline{\Iion}_{,\,c}
  \;=\; \frac{1}{|\Omegac|}\int_{\Omegac} \Iion \,\mathrm{d}\Omega
  \;\approx\; \frac{1}{|\Omegac|}\sum_{q} w_q\,\phi\!\big(\Vm(x_q),\,\states(x_q)\big).
  \label{eq:quad}
\end{equation}
The membrane potential at the Gauss points, $\Vm(x_q)$, is obtained in both modes
by the same high-order reconstruction of the cell-centred field. The two modes
differ \emph{only} in how the state vector $\states(x_q)$ that enters
\eqref{eq:quad} is obtained.

\paragraph{\gpode\ (legacy).}
The stiff ODE is integrated \emph{independently at every Gauss point}, driven by
the locally reconstructed potential:
\begin{equation}
  \states_q^{\,n+1}
  \;=\; \mathrm{ODEsolve}\!\big(\states_q^{\,n},\,\Vm(x_q);\,\Delta t\big),
  \qquad
  \states(x_q) \equiv \states_q^{\,n+1}.
  \label{eq:gp}
\end{equation}
Cost: $N_q$ stiff integrations per cell.

\paragraph{\ccrec\ (new).}
The ODE is integrated \emph{once per cell}, driven by the cell-centre potential
$\Vbarc \equiv \Vm(x_c)$, and the resulting cell-centred states are then
reconstructed to the Gauss points by an LRE of order $p$:
\begin{equation}
  \states_c^{\,n+1}
  \;=\; \mathrm{ODEsolve}\!\big(\states_c^{\,n},\,\Vbarc;\,\Delta t\big),
  \qquad
  \states(x_q) \equiv \widehat{\states}^{\,(p)}(x_q)
        = \mathcal{R}^{(p)}\!\big[\{\states_{c'}^{\,n+1}\}_{c'\in\text{stencil}}\big](x_q).
  \label{eq:cc}
\end{equation}
Cost: one stiff integration per cell (the motivation: $I_{\mathrm{Na}}$-limited
TNNP integrations dominate the run time, so $1$ vs $N_q$ ODEs per cell is a large
saving).

The essential contrast is already visible in \eqref{eq:gp}--\eqref{eq:cc}: in
\gpode\ the non-linear map $\Vm\!\mapsto\!\states$ is applied to the
\emph{pointwise} $\Vm(x_q)$; in \ccrec\ it is applied to the \emph{cell-average}
$\Vbarc$, and only the (already-computed) result is interpolated.

% ================================================================== %
\section{Observed symptom}
% ================================================================== %

Benchmark: 2-D slab, structured hexahedra, $\Delta x=\SI{0.5}{mm}$
($40\times16$), TNNP, a \SI{1.5}{mm} corner stimulus box of intensity
\SI{50000}{A/m^3} for \SI{2}{ms}, activation threshold \SI{0}{mV},
Crank--Nicolson, Picard coupling, RKF45 for the states. With \ccrec\ at order
\code{p1} (``CCp1''):
\begin{itemize}[nosep]
  \item at $\Delta t=\SI{e-4}{s}$ the activation ignites and propagates
        (diagonal points activate in sequence);
  \item at $\Delta t=\SI{e-5}{s}$ \emph{nothing} activates --- every diagonal
        sample, including the stimulus site itself, stays at
        $t_{\mathrm{act}}=0$.
\end{itemize}
A finer (more accurate) time step producing \emph{no} activation is the
counter-intuitive signature that must be explained.

% ================================================================== %
\section{Discrimination experiments}
% ================================================================== %

Three controlled runs isolate the cause. All use the same mesh
($40\times16$, $\Delta x=\SI{0.5}{mm}$), $\Delta t=\SI{e-5}{s}$, endtime
\SI{10}{ms} (long enough: the stimulus site fires at ${\sim}\SI{1.3}{ms}$ when it
triggers). ``Ignition'' means the stimulus site (diagonal index $0$) crosses the
threshold; ``conduction'' means the front advances along the diagonal.

\begin{table}[h]
\centering
\footnotesize
\begin{tabular}{@{}l l c c p{5.2cm}@{}}
\toprule
Run & Configuration & Ignition & Conduction & Evidence \\
\midrule
A & \gpode\ (\code{states=NO}), CN            & \textbf{yes} & \textbf{yes}
  & pt.\ $0$ @ \SI{1.20}{ms}; pts.\ $0$--$4$ active, front ${\sim}\SI{4.5}{mm}$ \\
--- & \ccrec\ CCp1, CN (original)             & no  & ---
  & all $t_{\mathrm{act}}=0$; $\max \Vm=\SI{-27}{mV}$ \\
B & \ccrec\ CCp1, \textbf{backwardEuler}      & no  & ---
  & all $t_{\mathrm{act}}=0$; $\max \Vm=\SI{-27}{mV}$ \\
C & \ccrec\ CCp1, CN, stimulus $\times 4$     & \textbf{yes} & no
  & pt.\ $0$ @ \SI{0.60}{ms}; pts.\ $0$--$2$ only, front stalls ${\sim}\SI{2.3}{mm}$ \\
\bottomrule
\end{tabular}
\caption{Discrimination runs at $\Delta t=\SI{e-5}{s}$. \gpode\ ignites and
conducts; \ccrec\ does neither with the nominal stimulus, and even a
$4\times$ stimulus only ignites without sustaining conduction.}
\label{tab:disc}
\end{table}

\noindent The three runs rule out the ``obvious'' suspects and localise the cause:
\begin{description}[style=nextline,leftmargin=1.2em,nosep]
  \item[Not the time scheme.] Run~B replaces Crank--Nicolson by the fully
    dissipative backward Euler and still fails identically
    ($\max\Vm=\SI{-27}{mV}$). Crank--Nicolson is exonerated.
  \item[Not a $\Delta t$-scaling bug.] The $\theta$-scheme is $\Delta t$-consistent
    ($\mathsf{A}=\mathsf{M}/\Delta t-\theta\mathsf{L}$,
    $\mathsf{B}=\mathsf{M}/\Delta t+(1-\theta)\mathsf{L}$, source weighted by
    $\theta,(1-\theta)$ with no spurious $\Delta t$ factor), and the stimulus is
    applied over a fixed physical window $[t_0,t_0+\SI{2}{ms}]$ independent of
    $\Delta t$. Both $\Delta t$ inject the same charge.
  \item[It \emph{is} the state mode.] Run~A is identical to the failing case
    except that the states are solved at the Gauss points (\gpode). It ignites at
    \SI{1.20}{ms} and conducts. The only changed ingredient is
    \eqref{eq:gp} vs \eqref{eq:cc}.
\end{description}

% ================================================================== %
\section{Root cause: averaging a non-linear kinetics}
% ================================================================== %

\subsection{The sodium activation is convex at the foot of the upstroke}

The fast depolarisation is carried by the sodium current
$I_{\mathrm{Na}}=g_{\mathrm{Na}}\,m^{3}h\,j\,(\Vm-E_{\mathrm{Na}})$, whose
activation gate relaxes towards the steep sigmoid
\begin{equation}
  m_\infty(\Vm)=\Big[1+\exp\!\big((V_{1/2}-\Vm)/k\big)\Big]^{-1},
  \qquad V_{1/2}\approx\SI{-40}{mV},\; k\approx\SI{6}{mV}.
\end{equation}
Below its inflection ($\Vm<V_{1/2}$) --- exactly the sub-threshold ``foot'' where
ignition is decided --- $m_\infty$ is \emph{convex}. Over a cell $\Omegac$ that
straddles the wavefront, $\Vm(x)$ ranges from resting on the trailing side to
strongly depolarised on the leading side.

\subsection{Jensen's inequality: \ccrec\ under-shoots the reactive current}

Let $\avg{\cdot}=\tfrac{1}{|\Omegac|}\sum_q w_q(\cdot)$ denote the cell-quadrature
average, so $\Vbarc=\avg{\Vm}$. \gpode\ evaluates the kinetics pointwise and
averages the \emph{outcome}; \ccrec\ averages the \emph{input} first. For the
convex $m_\infty$,
\begin{equation}
  \underbrace{\avg{\,m_\infty(\Vm)\,}}_{\text{\gpode\ (effective)}}
  \;\;\ge\;\;
  \underbrace{m_\infty\!\big(\avg{\Vm}\big)}_{\text{\ccrec}}
  \;=\; m_\infty(\Vbarc).
  \label{eq:jensen}
\end{equation}
Because $I_{\mathrm{Na}}\propto m^{3}$, the gap in \eqref{eq:jensen} is amplified
cubically. \ccrec\ therefore systematically \emph{under-estimates} the open-sodium
fraction, hence the inward (depolarising) current, precisely where $\Vm$ varies
most within a cell --- at the stimulus edge and along the upstroke. In \gpode\ the
Gauss points sitting on the depolarised side of the cell each run their own ODE,
see a high local $\Vm(x_q)$, open their sodium channels, and contribute a strong
inward current to \eqref{eq:quad}; those sub-cell ``hot spots'' are exactly what
\ccrec\ averages away before the kinetics ever sees them.

\subsection{Two failure modes}

\begin{enumerate}[leftmargin=1.4em,nosep]
  \item \textbf{Ignition failure.} With the nominal stimulus the site depolarises
    to only $\Vm\approx\SI{-27}{mV}$ (Runs~orig/B) and relaxes back: the weakened
    $I_{\mathrm{Na}}$ never becomes regenerative, so the all-or-none upstroke does
    not fire. A $4\times$ stimulus (Run~C) forces $\Vm$ past threshold at the site
    ($\SI{0.60}{ms}$), confirming the nominal stimulus was rendered
    \emph{effectively sub-threshold} by the averaging.
  \item \textbf{Conduction block.} Even once ignited (Run~C), the front stalls
    after $2$--$3$ cells. Propagation requires each activating cell to supply
    enough inward charge (\emph{source}) to bring the downstream cells
    (\emph{sink}) to threshold --- the cardiac ``safety factor''. By
    \eqref{eq:jensen}, \ccrec\ lowers the source at the very front where it is
    needed, driving the safety factor below unity: a numerically induced
    source--sink mismatch, i.e.\ conduction block.
\end{enumerate}

% ================================================================== %
\section{Why the coarse time step ``worked''}
% ================================================================== %

At $\Delta t=\SI{e-4}{s}$ the implicit step advances the depolarisation in one
large jump: the reaction/stimulus source acts over a big $\Delta t/\mathsf{M}$
increment, producing a transient over-shoot of $\Vbarc$ that can exceed threshold
even with the under-estimated $I_{\mathrm{Na}}$. As $\Delta t\to0$ the computed
trajectory converges to the \emph{true} (accurate) one, which under \ccrec's
weakened current stays sub-threshold --- so the deficiency, masked by
coarse-step over-shoot, is revealed. Consequently the $\Delta t=\SI{e-4}{s}$
CCp1 activation is an artefact of temporal under-resolution and its activation
times should not be trusted; it is worth checking that they agree with \gpode\ at
the same $\Delta t$.

% ================================================================== %
\section{Why higher reconstruction order does not help}
% ================================================================== %

Modes \code{p2} and \code{p3} change only the operator $\mathcal{R}^{(p)}$ in
\eqref{eq:cc} --- how the cell-centred states $\states_c^{\,n+1}$ are interpolated
to the Gauss points. They do \emph{not} change the fact that $\states_c^{\,n+1}$
was produced by integrating the ODE with the cell-average input $\Vbarc$. The
bias in \eqref{eq:jensen} is committed \emph{inside} the ODE solve, before any
reconstruction; a higher-order interpolation merely refines an already-biased
quantity. Empirically and by construction, CCp2/CCp3 inherit the same ignition
and conduction failure as CCp1 on this front.

% ================================================================== %
\section{When \ccrec\ is valid, and recommendation}
% ================================================================== %

\eqref{eq:jensen} is an equality when $\Vm$ is (nearly) linear across the cell,
i.e.\ when the sub-cell variation is small compared with the curvature scale of
$m_\infty$. This is the condition
\begin{equation}
  h_{\text{cell}} \;\ll\; w_{\text{front}},
\end{equation}
the cell size being much smaller than the upstroke width. The cardiac upstroke is
${\sim}\SI{0.5}{}$--$\SI{1}{mm}$ wide; at $h=\SI{0.5}{mm}$ a cell spans the whole
front and \ccrec\ fails, whereas $h\lesssim\SI{0.1}{mm}$ would resolve it --- at a
cell count that erodes the very cost advantage \ccrec\ was meant to provide.

\begin{itemize}[leftmargin=1.4em,nosep]
  \item \textbf{Propagation with a stiff cell model (Niederer/TNNP):} use
    \gpode\ (\code{stateIntegrationMode gaussPointODE}, or \code{states="NO"} in
    the triad sweep). It resolves the sub-cell kinetics and both ignites and
    conducts.
  \item \textbf{Smooth reaction fields} (e.g.\ the manufactured MMS case of the
    FDA solver, where $\Vm$ varies little within a cell): \ccrec\ is accurate and
    delivers its intended $N_q\!\to\!1$ ODE speed-up.
  \item \textbf{Possible hybrid (future work):} integrate the ODE at the Gauss
    points only where $\lvert\grad\Vm\rvert$ is large (the moving front) and at
    the cell centre with reconstruction elsewhere (resting/plateau). This would
    recover most of the speed-up while preserving accuracy at the front.
\end{itemize}

\paragraph{Bottom line.} \ccrec\ is not broken and its implementation is correct;
it simply trades sub-cell resolution of the ionic kinetics for speed. That trade
is safe for smooth reactions but invalid for the steep, self-sustaining upstroke
of a real cardiac action potential, where averaging the membrane potential before
the sodium kinetics suppresses ignition and blocks conduction. For the Niederer
benchmark, \gpode\ is the appropriate mode.

\end{document}
